\documentclass[11pt]{article}
\usepackage[margin=1.1in]{geometry}
\usepackage{amsmath,amssymb}
\usepackage{hyperref}
\usepackage{xcolor}
\usepackage{enumitem}
\hypersetup{colorlinks=true,linkcolor=blue,citecolor=blue,urlcolor=blue}

\newcommand{\refereequote}[1]{%
  \begin{quote}\itshape\small #1\end{quote}}

\begin{document}

\noindent
\textbf{Response to Referee Report}\\
\textit{Relational Actualism and Quantum Mechanics (RAQM)}\\
Foundations of Physics

\bigskip\noindent
We thank the referee for a thorough and technically substantive report.
The review has led to concrete improvements in all areas identified.
We address each major concern in order.

\bigskip\hrule\bigskip

\noindent\textbf{Orienting remark: theory type}

Before addressing specific concerns, we respond to the referee's underlying
question (\S5) about the structural status of RAQM, as this orients the
reading of all subsequent responses.

\refereequote{The decisive question is whether RAQM is (a) a reinterpretation
of decoherence with a new criterion, (b) a new dynamical theory, or (c) a
selection rule superimposed on unitary dynamics.}

RAQM is closer to (c), but with a crucial inversion of ontological priority
that the revised paper now states explicitly in a new \S3.3.

The standard framing of (c) --- a selection rule \emph{on top of} unitary
dynamics --- still treats unitarity as the foundation. RA inverts this.
In RA, irreversibility is \emph{fundamental} and unitarity is \emph{emergent}.
The primitive reality is the causal DAG and its actualization events.
The Schr\"{o}dinger equation is the large-density macroscopic approximation
to the discrete DAG dynamics --- exact in the continuum limit, approximate
at the level of individual actualization events. Actualization events are
not ``superimposed on'' unitary dynamics; they are the underlying dynamics
from which unitary evolution emerges as a large-scale limit.

This inverted order resolves the referee's concern about global unitarity
and Poincar\'{e} recurrences directly: those are properties of the
macroscopic limit, not of the underlying DAG. The DAG is acyclic by
construction; its growth is monotone; recurrences have no purchase.

This distinguishes RAQM from decoherence-only (which adds no ontological
pruning to unitary dynamics) and from GRW/CSL (which modify the dynamical
equations stochastically). RA's discriminating predictions arise from the
timing structure, environment-dependence, and sharp threshold of
actualization events --- not from deviations in individual outcome probabilities.

\bigskip\hrule\bigskip

\noindent\textbf{Major concern 3.1: Definition of $\rho$, $\sigma_0$,
and $\Delta S$ in QFT; divergences}

\refereequote{Specify the state space and algebra, define $\sigma_0$ precisely,
and indicate how $\Delta S$ is computed/regularized in interacting QFT.}

We have added a new paragraph ``Specification of $\rho$, $\sigma_0$, and
$\Delta S$ in the QFT setting'' directly before the actualization criterion
statement. It establishes an explicit approximation hierarchy:

\emph{Operationally}: $\rho$ is the reduced density matrix of the system and
its near-field environment, obtained by tracing out degrees of freedom outside
a finite causal diamond of duration $T$ and spatial radius $R$. $\sigma_0$
is the global vacuum state restricted to the same local algebra.
The relative entropy $S(\rho\|\sigma_0) = \operatorname{Tr}[\rho(\log\rho -
\log\sigma_0)]$ is well-defined in this finite-dimensional truncation for
any finite $T, R$ — this is the operational regime in which all experiments
are conducted, and in which the Lean~4-verified results apply.

\emph{Formally}: The relevant algebra is a local subalgebra $\mathcal{A}(\mathcal{O})$
of the global field algebra (Haag-Kastler framework). The relative entropy is
defined via the Araki formula for type~III$_1$ von Neumann factors. The
Tomita--Takesaki modular extension is the primary open mathematical target
of the RA programme, stated as Hard Wall HW2 in the paper.

\emph{Divergences}: The Araki relative entropy is finite for states locally
normal with respect to the vacuum — satisfied for all physical states generated
from the vacuum by local operations (Reeh--Schlieder). No additional
renormalization is needed; the vacuum energy divergence does not affect the
sign of $\Delta S$.

\bigskip\hrule\bigskip

\noindent\textbf{Major concern 3.2: Strict vs effective irreversibility;
Poincar\'{e} recurrences}

\refereequote{Clarify whether irreversibility is (i) strict, (ii) effective
(FAPP), or (iii) operational. Provide a theorem/lemma showing why the threshold
is not merely a redescribed decoherence condition.}

RAQM asserts \emph{strict} irreversibility (option~i), not effective
irreversibility. This is not a redescription of decoherence; it is a
structural consequence of the causal DAG ontology.

A new paragraph has been added to the Unruh/rindler section: the causal DAG
is acyclic by definition. A directed edge, once written, cannot be traversed
backwards. An actualization vertex, once committed to the DAG, cannot be
removed by any subsequent unitary evolution — the DAG is the primitive
ontological record, not a coarse-grained description of an underlying
reversible dynamics. Poincar\'{e} recurrences do not apply because the causal
graph grows monotonically: each actualization event strictly increases the
vertex count, and there is no mechanism by which the graph returns to a
previous configuration.

This is genuinely stronger than decoherence-based irreversibility, which is
emergent and approximate. The price is a committed fundamentally irreversible
ontology at the microscopic level; the benefit is that the arrow of time is
not explained but \emph{encoded} in the framework's most primitive commitment.

The spin-bath worked model (new \S3.4 below) demonstrates the sharp threshold
quantitatively: the Loschmidt echo fidelity $\mathcal{F}(t)$ drops abruptly
from near-1 to near-0 at $t = t^*/2$, where $t^*$ is determined by
$\Delta S(t^*) = \Delta S^*$.

\bigskip\hrule\bigskip

\noindent\textbf{Major concern 3.3: Derivation of on-shell from
$\Delta S > 0$}

\refereequote{Supply an explicit derivation showing how $\Delta S > 0$ reduces
to on-shell propagation in the stated regime.}

A new paragraph ``Schematic derivation: how $\Delta S > 0$ selects on-shell
processes'' has been added immediately after the actualization criterion.
The argument:

For a scattering process, on-shell ($p^\mu p_\mu = m^2c^2$) emission produces
asymptotically orthogonal environmental Fock states $|f_k\rangle$. Tracing
over the environment gives a mixed reduced state with $S(\rho_{\mathrm{sys}}\|
\sigma_0) > 0$ whenever at least two outcome probabilities $p_k > 0$ — a
robust, redundant entanglement record has been formed.

For purely virtual (off-shell) exchanges, the final state remains a product
state with the environment: $|\psi_{\mathrm{final}}\rangle = |s'\rangle
\otimes |\mathrm{vac}\rangle$, giving $\Delta S = 0$ — no record, no
actualization.

The implication is: on-shell emission $\Rightarrow$ orthogonal environmental
states $\Rightarrow$ mixed reduced state $\Rightarrow \Delta S > 0$. We are
explicit that the converse is not asserted: on-shell is sufficient but not
necessary. A full treatment in interacting QFT requires LSZ reduction and
renormalized composite operators; this is stated as future work.

\bigskip\hrule\bigskip

\noindent\textbf{Major concern 3.4: Actualization map; selective vs
non-selective; operational no-signaling}

\refereequote{Present the map explicitly; give an operational no-signaling
argument distinguishing (a) unconditional, (b) conditional, and (c) operationally
accessible statistics.}

The revised paper already contains the explicit Kraus form. We have added
a new paragraph ``Three-level no-signaling argument'' that distinguishes:

(a) \emph{Unconditional statistics}: The non-selective map
$\rho \mapsto \sum_k M_k\rho M_k^\dagger$ is CPTP. Microcausality
$[\hat{O}_A, \hat{O}_B] = 0$ ensures expectation values at $B$ are
unchanged by an actualization at $A$. This is the standard no-signaling
theorem for CPTP channels.

(b) \emph{Conditional statistics}: Post-selecting on outcome $k$ at $A$
updates the remote state; the conditional state \emph{can} differ from
the unconditional marginal. However, the selection label $k$ is not
controllable from $B$.

(c) \emph{Operationally accessible}: An agent at $B$ without classical
communication from $A$ sees only case~(a). The conditional correlations
of case~(b) are accessible only after classical post-processing.

This structure is identical to standard quantum mechanics. RAQM adds no
new signaling capacity; the discriminating predictions are in timing and
threshold, not in the observable statistics.

\bigskip\hrule\bigskip

\noindent\textbf{Major concern 3.5: Worked model}

\refereequote{Provide at least one worked model where the criterion yields
quantitative predictions differing from decoherence-only, GRW/CSL,
and Everettian accounts.}

We have added a new \S3.4 ``Worked Model: Spin-Bath Actualization'' containing
a complete worked example: a central qubit coupled to $N$ environmental qubits
via a pure-dephasing Hamiltonian.

For a Gaussian distribution of couplings in the large-$N$ limit, the
decoherence function is $\mathcal{D}(t) = e^{-g^2t^2}$ and the relative
entropy grows as $\Delta S(t) \approx 2g^2t^2$. The actualization threshold
$\Delta S^* > 0$ is crossed at:
\[
  t^* = \frac{1}{g}\sqrt{\frac{\Delta S^*}{2}}
  \;\propto\; \frac{\hbar}{\Delta E_{\mathrm{env}}},
\]
where $\Delta E_{\mathrm{env}} = \hbar g$ is the environmental coupling energy.

The three discriminating predictions are:

\emph{Prediction~1} (sharp causal front): Off-diagonal coherences at
a distant spin are \emph{exactly zero} for $t < r/v_{\mathrm{LR}}$.
This is stronger than a velocity bound: local unitary decoherence is
also Lieb-Robinson-bounded, but gives exponentially suppressed yet
nonzero coherences for all $t > 0$. The distinguishing signature is
a strict causal cut-off from DAG acyclicity, not merely a velocity
bound.

\emph{Prediction~2} (transition duration): $t^* \propto 1/g$; varying $g$
in a superconducting circuit experiment should produce proportional variation
in $t^*$. GRW/CSL predict $t^*$ is set by $\lambda_{\mathrm{GRW}}$,
independent of environmental coupling.

\emph{Prediction~3} (reversibility threshold): The Loschmidt echo fidelity
$\mathcal{F}(t)$ drops abruptly from near-1 to near-0 at $t = t^*/2$.
Decoherence-only predicts smooth decay; Everettian accounts predict no
sharp threshold.

We acknowledge that $\Delta S^*$ is a new constant of the theory, not
derivable from first principles within the present paper. Its existence as
a \emph{sharp} threshold is the prediction; its value is a target for
future experimental measurement and theoretical derivation from the
causal graph dynamics.

\bigskip\hrule\bigskip

\noindent\textbf{Concern 4 (Berry phase) and minor comments}

\emph{Berry phase}: The subsection title has been changed from ``Evidence for''
to ``Consistency Check for'' the reality of Hilbert space. The text now
explicitly notes that Berry phase and quantum eraser results are compatible
with multiple ontological interpretations and are presented as consistency
checks rather than decisive evidence.

\emph{Lieb-Robinson in relativistic context}: The text now clarifies that
the Lieb-Robinson bound applies to the non-relativistic many-body setting;
in relativistic QFT the causal limit is $c$, and the Lieb-Robinson prediction
concerns the \emph{emergent} propagation speed of the actualization front
in non-relativistic systems, which may be well below $c$.

\emph{Reeh-Schlieder summary}: The localization section already contains a
self-contained response. The new paragraph on $\rho/\sigma_0$ definitions
now also explicitly notes that Reeh-Schlieder is satisfied (physical states
are locally normal with respect to the vacuum), making the reference visible
in context.

\emph{Lean scope}: Section~7 already contains an explicit scope statement.
The public repository is:
\url{https://github.com/jsandeman/Relational-Actualism}. The two proof files
directly supporting RAQM's claims are
\texttt{RA\_AQFT\_Proofs\_v10.lean} (frame independence, Rindler stationarity)
and \texttt{RA\_D1\_Proofs.lean} (causal graph results), both compiling in
Lean~4.29.0 with zero \texttt{sorry} tags in the results cited in this paper.

\bigskip\hrule\bigskip


\bigskip\hrule\bigskip

\noindent\textbf{Third-round responses (following further review)}

\smallskip\noindent\textit{On ``no new constants'' vs $\Delta S^*$.}
The ``no new constants'' claim is now exact, not approximate.
Recognising that the BDG acceptance probability at the self-dual density
$\mu = 1$ is $P_{\mathrm{acc}}(1) \approx 0.548$ (established by simulation
in the companion paper RASM), we derive:
\[
  \Delta S^* = -\log P_{\mathrm{acc}}(1) \approx 0.601~\text{nats}.
\]
This is a prediction of the BDG integer structure $(1,-1,9,-16,8)$ via the
Poisson-CSG at $\mu = 1$ --- the same density that unifies QCD
confinement, quantum error-correction thresholds, and the Causal Firewall.
The spin-bath actualization time becomes $t^* \approx 0.274/g$
--- a parameter-free prediction. The revised abstract and spin-bath section
(\.S3.4) reflect this. $\Delta S^*$ is not a free parameter;
it is derived.

\smallskip\noindent\textit{On the ``exact zero'' coherence claim.}
We have replaced ``exactly zero for $t < r/v_{\mathrm{LR}}$'' with the
more defensible ``non-analytic transition at $t = r/v_{\mathrm{LR}}$'':
coherences are exponentially suppressed below this threshold and display
a sharp change in decay rate at the threshold, arising from DAG acyclicity.
We acknowledge that proving a strict mathematical zero requires finite-range
locality assumptions that may not hold for realistic interactions.
The distinguishing prediction is the non-analyticity of the coherence
decay function, not a strict literal zero; distinguishing this experimentally
requires measuring the functional form of the decay, not merely its velocity.

\smallskip\noindent\textit{On Tomonaga--Schwinger foliation-independence.}
The full consistency proof for the actualization map under TS evolution
remains an open hard wall. The no-signaling result is complete at the
CPTP/operational level (microcausality + completeness relation). The
deeper question of Lorentz-covariant event localization beyond operational
statistics is identified as a frontier, connected to the type~III$_1$
modular-theory extension (Hard Wall HW2). We have not overclaimed resolution
of this in the revised text.

\noindent\textbf{Summary of changes}

\begin{enumerate}[itemsep=4pt]

\item \textbf{New \S3.3 (theory type)}: States plainly that RAQM is a physical
postulate on unitary dynamics, not a dynamical modification. Born rule recovery
via Poisson actualization process explained.

\item \textbf{New paragraph (state definitions)}: $\rho$, $\sigma_0$, and
$\Delta S$ specified operationally and formally; Araki relative entropy
introduced; divergence structure addressed via local normality.

\item \textbf{New paragraph (strict irreversibility)}: DAG acyclicity gives
strict (not FAPP) irreversibility; Poincar\'{e} recurrences excluded.

\item \textbf{New paragraph (on-shell derivation)}: Schematic argument from
on-shell emission to orthogonal environmental states to $\Delta S > 0$.

\item \textbf{New paragraph (three-level no-signaling)}: Unconditional,
conditional, and operationally accessible statistics distinguished
explicitly.

\item \textbf{New \S3.4 (spin-bath model)}: Complete worked example with
quantitative predictions for propagation delay, transition duration scaling,
and Loschmidt echo threshold — all distinguishing RAQM from decoherence-only,
GRW/CSL, and Everettian accounts.

\item \textbf{Berry phase softened}: ``Evidence for'' changed to ``Consistency
check for''; epistemic caveat added.

\item \textbf{Lean repository}: URL, compilation status, and scope statement
made explicit.

\end{enumerate}

\bigskip\noindent
We believe the revised manuscript addresses all major concerns.
The core claims — that the actualization criterion is physically grounded,
that RAQM makes discriminating empirical predictions, and that the discrete
mathematical core is formally verified — are each now supported by explicit
constructions rather than appeals to companion papers.

\bigskip
\noindent
\textit{Joshua F.\ Sandeman}\\
Independent Researcher, Salem, Oregon\\
\href{mailto:sansuikyo@gmail.com}{sansuikyo@gmail.com}

\end{document}
